\begin{table}[H]
\caption{Чувствительность четырех основных контрастов к старым и новым задачам. Значения указаны в процентных пунктах с 95\%-ми интервалами бутстрепа по задачам.}
\label{tab:subgroup-sensitivity}
\centering\small
\begin{tabular}{@{}lrrr@{}}
\toprule
Контраст & Подмножество & Задачи & Разность [95\% интервал]\\
\midrule
SR--SN & повторные 200 & 194 & $-1{,}03\;[-3{,}78;\,1{,}72]$\\
       & новые 800 & 788 & $-0{,}68\;[-2{,}07;\,0{,}72]$\\
MR--MN & повторные 200 & 190 & $+1{,}93\;[-0{,}70;\,4{,}56]$\\
       & новые 800 & 774 & $+1{,}21\;[-0{,}13;\,2{,}58]$\\
MN--SN & повторные 200 & 192 & $-4{,}17\;[-7{,}64;\,-0{,}87]$\\
       & новые 800 & 776 & $-4{,}17\;[-5{,}84;\,-2{,}58]$\\
MR--SR & повторные 200 & 190 & $-1{,}23\;[-3{,}68;\,1{,}23]$\\
       & новые 800 & 781 & $-2{,}35\;[-3{,}88;\,-0{,}81]$\\
\bottomrule
\end{tabular}
\par\smallskip\noindent Каждый ответ был сгенерирован заново. «Повторные» означает, что идентификатор задачи встречался в прежнем исследовании 200 задач; сам результат повторно не использовался.
\end{table}
